\newglossaryentry{IPSA_g}{
    name={Institut Polytechnique des Sciences Avancées},
    description={
        Ecole d'ingénieurs aéronautiques et spatiaux privée française.
    },
}
\newglossaryentry{PMI_g}{
    name={Projet Master \acrshort{IPSA}},
    description={
        Projet de fin de cursus d'ingénieur à l'\acrshort{IPSA}, visant à concevoir, réaliser et tester un système innovant en
        lien avec         les domaines de l'aéronautique et du spatial.
    },
}
\newglossaryentry{CNES_g}{
    name={Centre National d'Études Spatiales},
    description={
        Agence spatiale française, responsable des programmes spatiaux civils en France.
    },
}
\newglossaryentry{C'Space}{
    name={C'Space},
    description={
        Campagne nationale de lancements de fusées expérimentales organisée par le \acrshort{CNES} et l'association
        \gls{Planète Sciences}. Elle est exculsivement réservée aux étudiants des écoles d'ingénieurs et universités françaises.
        Cette campagne se déroule chaque année dans le camp de Ger à Tarbes dans les Hautes-Pyrénées au début du mois de juillet.
    },
}
\newglossaryentry{APEX}{
    name={APEX},
    description={
        Système de mesure embarqué développé par l'association AéroIPSA pour la collecte, l'enregistrement et la télémesure des
        données d'avionique durant le vol de fusée expérimentale. (\cite{APEX})
    },
}
\newglossaryentry{Planète Sciences}{
    name={Planète Sciences},
    description={
        Association française de promotion des sciences et des techniques auprès des jeunes, organisatrice de la campagne
        \gls{C'Space} en collaboration avec le \acrshort{CNES}.
    },
}
\newglossaryentry{CFD_g}{
    name={Computational Fluid Dynamics},
    description={
        Mécanique des fluides numérique, méthode de simulation numérique des écoulements de fluides basée sur la résolution des
        équations de Navier-Stokes \cite{NavierStokes}.
    },
}
\newglossaryentry{FEM_g}{
    name={Finite Element Method},
    description={
        Méthode des éléments finis, technique de simulation numérique utilisée pour résoudre des problèmes d'ingénierie et de
        physique en divisant un domaine complexe en sous-domaines plus simples appelés éléments finis.
    }
}
\newglossaryentry{CAO_g}{
    name={Conception Assistée par Ordinateur},
    description={
        Ensemble de méthodes et d'outils informatiques permettant de concevoir, modéliser et simuler des objets ou systèmes
        techniques en 2D ou 3D.
    },
}
\newglossaryentry{AMDEC_g}{
    name={Analyse des Modes de Défaillance, de leurs Effets et de leur Criticité},
    description={
        Méthode systématique d'analyse des risques visant à identifier les modes de défaillance potentiels d'un système, à
        évaluer leurs effets et à déterminer leur criticité afin de mettre en place des actions correctives appropriées.
    },
}
\newglossaryentry{Fusex}{
    name={Fusex},
    description={
        Fusée expérimental : catégorie de fusée désigné par le \acrshort{CNES} et \gls{Planète Sciences} pour les fusées
        atmosphériques de grand calibre. 
    },
}
\newglossaryentry{Minif}{
    name={Minif},
    description={
        Mini fusée : catégorie de fusée désigné par le \acrshort{CNES} et \gls{Planète Sciences} pour les petites fusées
        atmosphériques de faible calibre.
    },
}
\newglossaryentry{Cansat}{
    name={Cansat},
    plural={Cansats},
    description={
        Petit satellite de la taille d'une canette de boisson, utilisé à des fins éducatives et expérimentales pour simuler
        certaines fonctions d'un satellite en orbite terrestre basse.
    },
}
\newglossaryentry{APR_g}{
    name={Analyse Préliminaire des Risques},
    description={
        Première étape de l'analyse de sûreté de fonctionnement visant à identifier les situations pouvant compromettre la
        sécurité du vol, l'intégrité du lanceur, ou la conformité aux exigences du \acrshort{CDC} \gls{Fusex} bi-étage.
    },
}
\newglossaryentry{MDF_g}{
    name={Modes de Défaillance},
    description={
        Causes techniques susceptibles de conduire aux \acrshort{ERPs} identifiés en \acrshort{APR}. Contrairement aux
        \acrshort{ERP}, qui décrivent des situations finales dangereuses au niveau système, les modes de défaillance décrivent
        les anomalies, ruptures ou dysfonctionnements pouvant apparaître au sein des sous-systèmes du lanceur. Les
        \acrshort{MDF} sont les points d'entrée de l'analyse détaillée de sûreté de fonctionnement, notamment dans le cadre de
        l'\acrshort{AMDEC}.
    },
}
\newglossaryentry{CATIA_g}{
    name={Conception Assistée Tridimensionnelle Interactive Appliquée},
    description={
        Logiciel de \acrshort{CAO} développé par Dassault Systèmes, utilisé pour la modélisation 3D, la simulation et la gestion
        du cycle de vie des produits dans divers domaines industriels, notamment l'aéronautique et l'automobile. C'est le
        logiciel enseigné à l'\acrshort{IPSA} et utilisé en majorité pour la conception des pièces du projet SP-02.
    },
}
\newglossaryentry{GNSS_g}{
    name={Global Navigation Satellite System},
    description={
        Système mondial de navigation par satellite : ensemble de constellations de satellites fournissant des services de
        positionnement, de navigation et de synchronisation temporelle à l'échelle mondiale. Les exemples les plus connus sont
        le GPS (États-Unis), Galileo (Union Européenne), GLONASS (Russie) et BeiDou (Chine).
    },
}
\newglossaryentry{LORA_g}{
    name={Long Range},
    description={
        Technologie de communication sans fil basse consommation et longue portée, utilisée pour la transmission de données sur
        de grandes distances avec une faible consommation d'énergie, souvent employée dans les applications de l'\acrfull{IoT}.
    },
}
\newglossaryentry{Skylab}{
    name={Skylab},
    description={
        Atelier de fabrication numérique de l'\acrshort{IPSA}, équipé de diverses machines-outils (imprimantes 3D, fraiseuses
        \acrshort{CNC}, découpeuses laser, etc.) permettant aux étudiants de prototyper et de fabriquer des pièces pour leurs
        projets techniques.
    },
}
\newglossaryentry{Stabtraj}{
    name={Stabtraj},
    description={
        Outil de dimensionnement et de simulation de trajectoire de fusée développé par l'association \gls{Planète Sciences} en
        collaboration avec le \acrshort{CNES}. Il permet d'analyser la stabilité, la performance et la trajectoire des fusées
        amateur. \cite{Stabtraj}
    },
}
\newglossaryentry{OpenRocket}{
    name={OpenRocket},
    description={
        Logiciel open-source de conception et de simulation de fusées amateur, permettant aux utilisateurs de modéliser des
        fusées, d'analyser leur stabilité et de simuler leurs trajectoires de vol. Il est largement utilisé par les passionnés
        d'astromodélisme.
    },
}
\newglossaryentry{RCE_g}{
    name={Rencontre Club Espace},
    description={
        Rencontre trimestrielle organisée par l'association \gls{Planète Sciences} et le \acrshort{CNES} pour les équipes
        participantes à la campagne \gls{C'Space}. Ces réunions permettent de faire le point sur l'avancement des projets,
        d'aborder les aspects techniques, réglementaires et logistiques, et de préparer les différentes étapes de la campagne de
        lancement.
    },
}
\newglossaryentry{IPR_g}{
    name={Indice de Priorité de Risque},
    description={
        Indicateur quantitatif utilisé dans l'analyse de sûreté de fonctionnement pour évaluer la criticité d'un mode de
        défaillance en fonction de la gravité de ses conséquences, de la fréquence d'apparition et de la détectabilité.
        L'\acrshort{IPR} permet de prioriser les actions correctives à mettre en place pour réduire les risques identifiés.
    },
}
\newglossaryentry{FAR_g}{
    name={Friends of Amateur Rocketry},
    description={
        Site de lancement de fusées expérimentales situé en Californie, aux États-Unis, utilisé par des amateurs et des
        professionnels pour tester et lancer des fusées de différents calibres.
    },
}
\newglossaryentry{ET_g}{
    name={Étude Technique},
    description={
        Désignation académique de sujet de \acrshort{PMI} à l'\acrshort{IPSA}. Chaque \acrshort{ET} porte un numéro unique et
        un intitulé décrivant la problématique technique abordée.
    },
}


\newacronym{IPSA}{IPSA}{\gls{IPSA_g}}
\newacronym{PMI}{PMI}{\gls{PMI_g}}
\newacronym{CNES}{CNES}{\gls{CNES_g}}
\newacronym{CFD}{CFD}{\gls{CFD_g}}
\newacronym{FEM}{FEM}{\gls{FEM_g}}
\newacronym{CAO}{CAO}{\gls{CAO_g}}
\newacronym{AMDEC}{AMDEC}{\gls{AMDEC_g}}
\newacronym{CDC}{CDC}{Cahier des Charges}
\newacronym{ERP}{ERP}{Événement Redouté Principal}
\newacronym{ERPs}{ERPs}{Événements Redoutés Principaux}
\newacronym{APR}{APR}{\gls{APR_g}}
\newacronym{MDF}{MDF}{\gls{MDF_g}}
\newacronym{CATIA}{CATIA}{\gls{CATIA_g}}
\newacronym{GNSS}{GNSS}{\gls{GNSS_g}}
\newacronym{LORA}{LORA}{\gls{LORA_g}}
\newacronym{EPI}{EPI}{Équipements de Protection Individuelle}
\newacronym{POMc}{POMc}{Polyoxyméthylène Copolymère}
\newacronym{CdG}{CdG}{Centre de Gravité}
\newacronym{CdP}{CdP}{Centre de Portance}
\newacronym{RCE}{RCE}{\gls{RCE_g}}
\newacronym{IoT}{IoT}{Internet des Objets (Internet of Things)}
\newacronym{CNC}{CNC}{Computer Numerical Control (Commande Numérique par Ordinateur)}
\newacronym{IPR}{IPR}{\gls{IPR_g}}
\newacronym{PLA}{PLA}{Polylactic Acid (Acide Polylactique)}
\newacronym{FAR}{FAR}{\gls{FAR_g}}
\newacronym{SNR}{SNR}{Signal-to-Noise Ratio (Rapport Signal sur Bruit)}
\newacronym{FSK}{FSK}{Frequency-Shift Keying (Modulation par Déplacement de Fréquence)}
\newacronym{ET}{ET}{\gls{ET_g}}
